% Transcription of Cayley's Collected Mathematical Papers, Vol.\ XII
% PDF pages 151--175 of collmathpapers12caylrich.pdf
% (book pages 131--155: Paper 825 -- Memoir on the Abelian and Theta Functions,
%  Chs.\ II conclusion, III in full, IV opening)
% Source scan: local workspace\Documents\Papors\OS\Cayley\collmathpapers12caylrich.pdf
% Restoration from scan; preserves Cayley's barred-determinant notation
% (\overline{012} etc.), Jacobian J(\phi,f), polar-function powers
% n.1^{n-1}2, and the symbols \partial_1, \partial_2 throughout.

\documentclass[11pt,a4paper]{article}

% ===== PACKAGES =====
\usepackage[utf8]{inputenc}
\usepackage[T1]{fontenc}
\usepackage[english]{babel}
\usepackage{amsmath,amssymb,amsthm}
\usepackage{geometry}
\usepackage{graphicx}
\usepackage{hyperref}
\usepackage{fancyhdr}
\usepackage{titlesec}
\usepackage{enumitem}
\usepackage{array}
\usepackage{booktabs}
\usepackage{longtable}
\usepackage{float}
\usepackage{caption}
\usepackage{subcaption}
\usepackage{mathrsfs}
\usepackage{mathtools}
\usepackage{bm}

% ===== GEOMETRY =====
\geometry{
  a4paper,
  margin=1in,
  top=1.2in,
  bottom=1in,
  headheight=14pt
}

% ===== HEADER/FOOTER =====
\pagestyle{fancy}
\fancyhf{}
\fancyhead[L]{\small Cayley, Collected Mathematical Papers, Vol.\ XII}
\fancyhead[R]{\small\thepage}
\renewcommand{\headrulewidth}{0.4pt}

% ===== TOLERANCES =====
\hfuzz=4pt
\tolerance=2000
\emergencystretch=10pt

% ===== THEOREM ENVIRONMENTS =====
\theoremstyle{plain}
\newtheorem{theorem}{Theorem}
\newtheorem{lemma}{Lemma}
\newtheorem{corollary}{Corollary}
\theoremstyle{definition}
\newtheorem{definition}{Definition}
\theoremstyle{remark}
\newtheorem{remark}{Remark}

% ===== CUSTOM COMMANDS =====
\providecommand{\paper}[2]{%
  \newpage
  \begin{center}
    {\Large\bfseries #1}\\[0.5em]
    #2
  \end{center}
  \vspace{1em}
}

\providecommand{\papertitle}[1]{%
  \begin{center}
    {\Large\scshape #1}
  \end{center}
  \vspace{0.5em}
}

\providecommand{\papersource}[1]{%
  \begin{center}
    {\small [#1]}
  \end{center}
  \vspace{0.5em}
}

\providecommand{\bookpage}[1]{\marginpar{\small\itshape p.~#1}}

% Elliptic function operators (modulus k, NOT i')
\DeclareMathOperator{\sn}{sn}
\DeclareMathOperator{\cn}{cn}
\DeclareMathOperator{\dn}{dn}

% Cayley's barred determinant
\providecommand{\bdet}[1]{\overline{#1}}

% ===== DOCUMENT =====
\begin{document}

\title{\scshape The Collected Mathematical Papers of Arthur Cayley\\Vol.\ XII\\[0.5em]\large Scan PDF Pages 151--175 (Book pp.\ 131--155; Paper 825, Chs.\ II concl.--IV)}
\author{Arthur Cayley (1821--1895)}
\date{Published 1889}
\maketitle

\thispagestyle{empty}
\begin{center}
\textit{Transcribed from the original Cambridge University Press edition, 1889.}\\
\textit{This file covers scan PDF pages 151--175 (book pages 131--155); restored from scan.}
\end{center}

\tableofcontents
\newpage

% ========================================================================
% PDF p. 151 = book p. 131 -- Paper 825 (cont.), Chapter II concluded
% ========================================================================
\section*{825. A Memoir on the Abelian and Theta Functions (continued)}
\addcontentsline{toc}{section}{825. A Memoir on the Abelian and Theta Functions (cont.)}

\subsection*{Chapter II. Proof of Abel's Theorem (concluded).}
\addcontentsline{toc}{subsection}{\quad Chapter II. Proof of Abel's Theorem (concl.)}

\noindent\textit{[book p.~131]}\quad \textbf{41.}\quad But for $\tau = 0$, we have
\[
A_x\phi_1 + B_x f_x = a\rho^{mn},
\]
\[
C_x\phi_1 + D_x f_x = \beta\sigma^{mn},
\]
and hence
\[
(AD - BC)_x f_x = A_x\beta\sigma^{mn} - C_x a\rho^{mn},
\]
and the right-hand side thus becomes, say
\[
-\Delta!\,\text{const.\ of}\;\left(-\,\frac{A_x}{a\rho^{mn}} + \frac{C_x}{\beta\sigma^{mn}}\right)\delta\phi_x.
\]

But in calculating the constant of $\dfrac{A_x}{a\rho^{mn}}\,\delta\phi_x$, we may suppose not only $\tau = 0$, but also $\sigma = 0$: we then have $\phi_x = (x,\,y,\,z)^m$, $= \left(\dfrac{\rho}{\Delta}\right)^m (x_1,\,y_1,\,z_1)^m$, $= \left(\dfrac{\rho}{\Delta}\right)^m \phi_1$, and hence also
\[
\delta\phi_x = \left(\frac{\rho}{\Delta}\right)^m \delta\phi_1.
\]

Similarly, in calculating the constant of $\dfrac{B_x}{\beta\sigma^{mn}}\,\delta\phi_x$, we may suppose not only $\tau = 0$, but also $\rho = 0$: we then have $\phi_x = (x,\,y,\,z)^m$, $= \left(\dfrac{\sigma}{\Delta}\right)^m (x_2,\,y_2,\,z_2)^m$, $= \left(\dfrac{\sigma}{\Delta}\right)^m \phi_2$, and hence
\[
\delta\phi_x = \left(\frac{\sigma}{\Delta}\right)^m \delta\phi_2.
\]

Moreover, in the equations
\[
A_x\phi_1 + B_x f_x = a\rho^{mn},
\]
\[
C_x\phi_1 + D_x f_x = \beta\sigma^{mn},
\]
writing in the first equation $\sigma = 0$, we find $A_x\left(\dfrac{\rho}{\Delta}\right)^m \phi_1 = a\rho^{mn}$; that is, $\dfrac{A_x}{a\rho^{mn}} = \left(\dfrac{\Delta}{\rho}\right)^m \dfrac{1}{\phi_1}$; and similarly writing in the second equation $\rho = 0$, we find $C_x\left(\dfrac{\sigma}{\Delta}\right)^m \phi_2 = \beta\sigma^{mn}$, that is, $\dfrac{C_x}{\beta\sigma^{mn}} = \left(\dfrac{\Delta}{\sigma}\right)^m \dfrac{1}{\phi_2}$; and the expression thus becomes
\[
= -\Delta!\,\left(-\,\frac{\delta\phi_1}{\phi_1} + \frac{\delta\phi_2}{\phi_2}\right),
\]
giving the equation which was to be proved. This completes the proof of the affected theorem
\[
\Sigma\,\frac{(x,\,y,\,z)_2^{n-2}\,d\omega}{012} = -\,\frac{\delta\phi_1}{\phi_1} + \frac{\delta\phi_2}{\phi_2}.
\]

\hfill $17$--$2$

% ========================================================================
% PDF p. 152 = book p. 132 -- Chapter III. The Major Function
% ========================================================================
\subsection*{Chapter III. The Major Function $(x,\,y,\,z)_2^{n-3}$.}
\addcontentsline{toc}{subsection}{\quad Chapter III. The Major Function}

\medskip\noindent\textit{Analytical Expression of the Function.}\quad Art.\ Nos.\ 42 to 49.

\medskip\noindent\textit{[book p.~132]}\quad \textbf{42.}\quad The function has been defined by the conditions that the curve $(x,\,y,\,z)_2^{n-3} = 0$ shall pass through the dps, and also through the $n - 2$ residues of the parametric points $1$, $2$; and moreover, that on writing therein $(x_1,\,y_1,\,z_1)$ for $(x,\,y,\,z)$, the function shall become $= n\,.\,1^{n-2}\,2$. Obviously the function is not completely determined: calling it $\Omega$ (or when required $\Omega_{12}$), then if $\Omega'$ be any particular form of it, the general form is $\Omega = \Omega' + (x,\,y,\,z)^{n-3}\,012$, where $(x,\,y,\,z)^{n-3} = 0$ is the general minor function, viz.\ $(x,\,y,\,z)^{n-3} = 0$ is a curve passing through the dps: the major function thus contains $\tfrac{1}{2}(n - 1)(n - 2) - \delta$, $= p$, arbitrary constants.

Agreeing with the definition, we have the before-mentioned equation
\[
\Omega = \frac{1}{\Delta^{n-1}}(n\,.\,1^{n-1}\,2,\,\ldots\!\,\Bigl)\Bigl(\rho,\,\sigma\Bigr)^{n-2} + \text{terms involving }\tau,
\]
viz.\ from this expression for $\Omega$ it appears that the curve $\Omega = 0$ meets the line through $1$, $2$ in the $n - 2$ residues of these points, and moreover, for $(x,\,y,\,z) = (x_1,\,y_1,\,z_1)$ and therefore $(\rho,\,\sigma,\,\tau) = (\Delta,\,0,\,0)$, the value of $\Omega$ is $= n\,.\,1^{n-2}\,2$.

\medskip\noindent\textbf{43.}\quad We can without difficulty write down an equation determining $\Omega'$ as a function $(x,\,y,\,z)^{n-3}$, which, on putting therein $\tau = 0$, becomes equal to the foregoing expression $\dfrac{1}{\Delta^{n-1}}(n\,.\,1^{n-1}\,2,\,\ldots\!\,\Bigl)\Bigl(\rho,\,\sigma\Bigr)^{n-2}$, and which is moreover such that the curve $\Omega' = 0$ passes through the dps; which being so, we have as before, $\Omega = \Omega' + (x,\,y,\,z)^{n-3}\,012$, for the general value of $\Omega$.

To fix the ideas, consider the particular case $n = 4$, the fixed curve a quartic: $\Omega'$, on putting therein $\tau = 0$, should become
\[
= \tfrac{1}{\Delta}(4\,.\,1^3\,2,\;6\,.\,1^2\,2^2,\;4\,.\,1\,2^{3}\!\,\Bigl)\Bigl(\rho,\,\sigma\Bigr)^{2};
\]
and it is be shown that this will be the case if we determine $\Omega'$, a quadric function of $(x,\,y,\,z)$, by the equation
\[
\left|\;\begin{array}{cc}
(x,\,y,\,z)^{2} & ,\quad \Omega'\\
1\,(x_1,\,y_1,\,z_1)^{2} & ,\quad 4\,.\,1^{3}\,2\\
2\,(x_1,\,y_1,\,z_1\;\xi\,x_2,\,y_2,\,z_2),\;6\,.\,1^{2}\,2^{2}\\
1\,(x_2,\,y_2,\,z_2)^{2} & ,\quad 4\,.\,1\,2^{3}\\
a,\;b,\;c,\;f,\;g,\;h, & 0
\end{array}\;\right| = 0,
\]
where the left-hand side is a determinant of seven lines and columns, the top line being $x^2$, $y^2$, $z^2$, $2yz$, $2zx$, $2xy$, $\Omega'$, and similarly for the second line; the third line is

% ========================================================================
% PDF p. 153 = book p. 133
% ========================================================================
\medskip\noindent\textit{[book p.~133]}\quad $2x_1x_2$, $2y_1y_2$, $2z_1z_2$, $2(y_1z_2 + y_2z_1)$, $2(z_1x_2 + z_2x_1)$, $2(x_1y_2 + x_2y_1)$, $6\,.\,1^{2}\,2^{2}$; and in each of the last three lines we have six arbitrary constants followed by a $0$. The equation is of the form $\Omega + M\Omega' = 0$, where $\Omega$ is a quadric function $(x,\,y,\,z)^2$, and $M$ is a constant factor.

\medskip\noindent\textbf{44.}\quad If the quartic curve has a dp, suppose at the point $a$, coordinates $(x_a,\,y_a,\,z_a)$, then in order that the curve $\Omega' = 0$ may pass through the dp, we must for one of the last three lines substitute $(x_a,\,y_a,\,z_a)^2$, $0$; and so for any other dp or dps of the quartic curve. And the conditions as to the dp or dps (if any) being satisfied in this manner, we may if we please, taking $(x_\beta,\,y_\beta,\,z_\beta)$ as the coordinates of an arbitrary point $\beta$ (not of necessity on the fixed curve), write any line not already so expressed, of the last three lines, in the form $(x_\beta,\,y_\beta,\,z_\beta)^2$, $0$; the effect being to make the curve $\Omega' = 0$ pass through the arbitrary point $\beta$.

\medskip\noindent\textbf{45.}\quad To show that the equation on putting therein $\tau = 0$ does in fact give the required value, $\Omega' = \dfrac{1}{\Delta}(4\,.\,1^3\,2,\;6\,.\,1^2\,2^2,\;4\,.\,1\,2^{3}\,\Bigl)\Bigl(\rho,\,\sigma\Bigr)^2 = \Phi$ suppose, it is to be observed that, effecting a linear substitution upon the first six columns, the equation may be written
\[
\left|\;\begin{array}{cc}
(\rho,\,\sigma,\,\tau)^{2}     & ,\quad \Omega'\\
1\,(\rho_1,\,\sigma_1,\,\tau_1)^{2} & ,\quad 4\,.\,1^{3}\,2\\
2\,(\rho_1,\,\sigma_1,\,\tau_1\;\xi\,\rho_2,\,\sigma_2,\,\tau_2),\;6\,.\,1^{2}\,2^{2}\\
1\,(\rho_2,\,\sigma_2,\,\tau_2)^{2} & ,\quad 4\,.\,1\,2^{3}\\
a',\;b',\;c',\;f',\;g',\;h', & 0
\end{array}\;\right| = 0,
\]
where $(\rho_1,\,\sigma_1,\,\tau_1)$, $(\rho_2,\,\sigma_2,\,\tau_2)$ are what $(\rho,\,\sigma)$ become on writing therein for $(x,\,y,\,z)$ the values $(x_1,\,y_1,\,z_1)$ and $(x_2,\,y_2,\,z_2)$ respectively, viz.\ we have $(\rho_1,\,\sigma_1,\,\tau_1) = (\Delta,\,0,\,0)$; $(\rho_2,\,\sigma_2,\,\tau_2) = (0,\,\Delta,\,0)$; the equation thus is
\[
\left|\;\begin{array}{ccccccc}
\rho^{2},  & \sigma^{2}, & \tau^{2}, & 2\sigma\tau, & 2\tau\rho, & 2\rho\sigma, & \Omega'\\
\Delta^{2}, & 0,         & 0,        & 0,           & 0,         & 0,           & 4\,.\,1^{3}\,2\\
0,         & 0,          & 0,        & 0,           & 0,         & 2\Delta^{2}, & 6\,.\,1^{2}\,2^{2}\\
0,         & \Delta^{2}, & 0,        & 0,           & 0,         & 0,           & 4\,.\,1\,2^{3}\\
a',        & b',         & c',       & f',          & g',        & h',          & 0
\end{array}\;\right| = 0,
\]
and then by another linear substitution upon the columns, the last column can be changed into $\Omega' - \Phi$, $0$, $0$, $0$, $0$, $0$, whence writing $\tau = 0$, the equation becomes
\[
\left|\;\begin{array}{ccccccc}
\rho^{2},  & \sigma^{2}, & 0, & 0, & 0, & 2\rho\sigma, & \Omega' - \Phi\\
\Delta^{2}, & 0,         & 0, & 0, & 0, & 0,           & 0\\
0,         & 0,          & 0, & 0, & 0, & 2\Delta^{2}, & 0\\
0,         & \Delta^{2}, & 0, & 0, & 0, & 0,           & 0\\
a',        & b',         & c',& f',& g',& h',          & 0
\end{array}\;\right| = 0,
\]

% ========================================================================
% PDF p. 154 = book p. 134
% ========================================================================
\medskip\noindent\textit{[book p.~134]}\quad or, omitting a constant factor, it is
\[
\left|\;\begin{array}{cccc}
\rho^{2}, & \sigma^{2}, & 2\rho\sigma, & \Omega' - \Phi\\
\Delta^{2}, & 0, & 0, & 0\\
0, & 0, & 2\Delta^{2}, & 0\\
0, & \Delta^{2}, & 0, & 0
\end{array}\;\right| = 0,
\]
that is, $\Omega' - \Phi = 0$, or $\Omega' = \Phi$, $= \dfrac{1}{\Delta}(4\,.\,1^3\,2,\;6\,.\,1^2\,2^2,\;4\,.\,1\,2^{3}\,\Bigl)\Bigl(\rho,\,\sigma\Bigr)^2$, the required value.

\medskip\noindent\textbf{46.}\quad Considering the equation for $\Omega'$ as expressed in the before-mentioned form $\Omega + M\Omega' = 0$, the value of the constant factor $M$ is
\[
M = \left|\;\begin{array}{c}
(x_1,\,y_1,\,z_1)^{2}\\
(x_1,\,y_1,\,z_1\;\xi\,x_2,\,y_2,\,z_2)\\
(x_2,\,y_2,\,z_2)^{2}\\
a,\;b,\;c,\;f,\;g,\;h,
\end{array}\;\right|,
\]
or if, instead of each line such as $a$, $b$, $c$, $f$, $g$, $h$, we have a line $(x_a,\,y_a,\,z_a)^2$, then we have
\[
M = \left|\;\begin{array}{c}
(x_1,\,y_1,\,z_1)^{2}\\
(x_1,\,y_1,\,z_1\;\xi\,x_2,\,y_2,\,z_2)\\
(x_2,\,y_2,\,z_2)^{2}\\
(x_3,\,y_3,\,z_3)^{2}\\
(x_4,\,y_4,\,z_4)^{2}
\end{array}\;\right|,
\]
a value which is
\[
= \;\left|\begin{array}{ccc} x_1, & y_1, & z_1\\ x_3, & y_3, & z_3\\ x_4, & y_4, & z_4 \end{array}\right|\cdot\left|\begin{array}{ccc} x_2, & y_2, & z_2\\ x_3, & y_3, & z_3\\ x_4, & y_4, & z_4 \end{array}\right|\cdot\left|\begin{array}{ccc} x_1, & y_1, & z_1\\ x_2, & y_2, & z_2\\ x_3, & y_3, & z_3 \end{array}\right|\cdot\left|\begin{array}{ccc} x_1, & y_1, & z_1\\ x_2, & y_2, & z_2\\ x_4, & y_4, & z_4 \end{array}\right|,
\]
or say this is $= 12s\,.\,12\beta\,.\,12\gamma\,.\,a\beta\gamma$.

\medskip\noindent\textbf{47.}\quad It is obvious that the foregoing process is applicable to the general case of the fixed curve of the order $n$ with $\delta$ dps, and gives always $\Omega'$, by an equation of the foregoing form $\Omega + M\Omega' = 0$, where $\Omega$ is a function $(x,\,y,\,z)^{n-3}$ of the coordinates, and $M$ is a constant factor. Supposing that in the determinant for $\Omega'$, each of the lower lines is written in the form $(x_a,\,y_a,\,z_a)^{n-3}$, $0$ instead of $a$, $b$, $c$, \&c., will indicate this function to be $= 0$; and for the points $\alpha$, are on a curve of the order $n - 3$.

% ========================================================================
% PDF p. 155 = book p. 135
% ========================================================================
\medskip\noindent\textit{[book p.~135]}\quad \textbf{48.}\quad Preceding the case $n = 4$, above considered, we have, of course, the case $n = 3$, $\delta = 0$, the fixed curve a cubic; the equation for $\Omega'$ is here
\[
\left|\;\begin{array}{cccc}
x, & y, & z, & \Omega'\\
x_1, & y_1, & z_1, & 3\,.\,1^{2}\,2\\
x_2, & y_2, & z_2, & 3\,.\,1\,2^{2}\\
x_3, & y_3, & z_3,
\end{array}\;\right| = 0,
\]
giving
\[
\tfrac{1}{3}\Omega' = \frac{1^{2}\,.\,023 + 1^{2}\,.\,031}{12s},
\]
or if we write herein $3$ for $a$, this is
\[
\tfrac{1}{3}\Omega' = \frac{1^{2}\,.\,023 + 1^{2}\,.\,031}{123},
\]
and we have hence the general form
\[
\tfrac{1}{3}\Omega = \frac{1^{2}\,.\,023 + 1^{2}\,.\,031}{123} + K\,.\,012,
\]
where $K$ is an arbitrary constant.

\medskip\noindent\textbf{49.}\quad There is, however, a more simple particular solution $\tfrac{1}{3}\Omega = $ polar function $012$ ($f = x^3 + y^3 + z^3$; then $012 = xx_1x_2 + yy_1y_2 + zz_1z_2$), to avoid a confusion of notation, we may write $= \bdet{012}$. We at once verify this; for, expressing the coordinates $(x,\,y,\,z)$ in terms of $(\rho,\,\sigma,\,\tau)$, we have $\tfrac{1}{3}\Omega = \bdet{012}$, $= \dfrac{1}{\Delta}(1^2\,.\,\rho + 1^2\,.\,\sigma + \bdet{125}\,.\,\tau)$, which, for $\tau = 0$ becomes $= \dfrac{1}{\Delta}\bigl[1^{2}\,.\,\rho + 1^{2}\,.\,\sigma\bigr]$.

We must, of course, have an identity of the form
\[
\bdet{012} = \frac{1^{2}\,.\,023 + 1^{2}\,.\,031}{123} + K\,.\,012,
\]
and to find $K$, writing here $0 = 3$, we have $K = \dfrac{\bdet{123}}{123^2}$, or we have the identity
\[
123\,\bdet{012} - \bdet{123}\,012 = 1^{2}\,.\,023 + 1^{2}\,.\,031.
\]

\medskip\noindent\textit{Single Letter Notation for the Polar Functions of the Cubic.}\quad Art.\ Nos.\ 50 and 51.

\medskip\noindent\textbf{50.}\quad The notation of single letters for the polar functions is not much required in the case of the cubic, but, in the next following case of the quartic it can hardly be dispensed with, and I therefore establish it in the case of the cubic: viz.\ I write
\[
2\bar{3},\;3\bar{1},\;1\bar{2} = f,\;g,\;h;\quad 23\bar{1},\;31\bar{2},\;12\bar{3} = i,\;j,\;k;\quad \bdet{123} = l,
\]
or, what is the same thing, the expression for the cubic function $f$ in terms of $\rho$, $\sigma$, $\tau$ is
\[
\Delta^2\,.\,f = 3l\rho^2\sigma + 3jp^2\tau + 6l\rho\sigma\tau + 6lp\sigma\tau + 3g\sigma^2\tau + 3i\sigma\tau^2 + 3i\sigma\tau^2;
\]

% ========================================================================
% PDF p. 156 = book p. 136
% ========================================================================
\medskip\noindent\textit{[book p.~136]}\quad an equation which, writing $0^3$ instead of $f$, may also be written
\[
\Delta!\,0^3 = (3l,\;3j,\;3i,\;3k,\;6l,\;6l,\;3i,\;3i,\;3j,\;3i\!\,\Bigl)\Bigl(\rho^3,\,\sigma^3,\,\tau^3,\,\rho^2\sigma,\,\rho^2\tau,\,\rho\sigma^2,\,\rho\tau^2,\,\sigma^2\tau,\,\sigma\tau^2\Bigr),
\]
and I join to it the series of equations
\begin{align*}
\Delta!\,.\,0^2\,1 &= (0,\;2k,\;2j,\;k,\;2l,\;2l,\;2l,\;g\!\,\Bigl)\,(\rho\sigma,\,\rho\tau,\,\sigma\tau,\ldots),\\
,,\;0^2\,2 &= (k,\;2l,\;2i,\;2l,\;2k,\;l\!\,\Bigl)\, ,\\
,,\;0^2\,3 &= (j,\;2l,\;2g,\;l,\;2h,\;2l\!\,\Bigl)\, ,\\
\Delta\,.\,0\,1^2 &= (0,\;h,\;l\!\,\xi\,\rho,\,\sigma,\,\tau),\\
\Delta\,\bdet{012} &= (h,\;l,\;l\!\,\xi\,\ldots\,),\\
,,\;0^2 &= (k,\;0,\;f\!\,\xi\;\,\ldots\,),\\
,,\;\bdet{025} &= (l,\;f,\;l\!\,\xi\;\,\ldots\,),\\
,,\;0\,3^2 &= (j,\;i,\;0\!\,\xi\;\,\ldots\,),
\end{align*}

\medskip\noindent\textbf{51.}\quad In particular, we have $\Delta\,\bdet{012} = h\rho + l\sigma + l\tau$, and the above-mentioned identity $123\,\bdet{012} - \bdet{123}\,012 = 1^2\,.\,023 + 1^2\,.\,031$ is simply $h\rho + l\sigma + l\tau - l\bdet{l} = h\rho + l\sigma$.

\medskip\noindent\textit{Single Letter Notation for the Polar Functions of the Quartic.}\quad Art.\ No.\ 52.

\medskip\noindent\textbf{52.}\quad I write here
\[
2\bar{3},\;3\bar{1},\;1\bar{2} = f,\;g,\;h;\quad 23\bar{1},\;31\bar{2},\;12\bar{3} = i,\;j,\;k;
\]
\[
1\bar{2}3,\;1\bar{2}3,\;1\bar{2}3 = l,\;m,\;n;\quad 23,\;31,\;1\bar{2} = p,\;q,\;r;
\]
so that the expression for the quartic function $f$ in terms of $\rho$, $\sigma$, $\tau$ is
\begin{align*}
\Delta!\,.\,f &= 4l\rho^3\sigma + 4jp^3\tau + 12lp^2\sigma^2 + 6gp^2\sigma^2\\
&\quad + 4lp\sigma + 12mp\sigma\tau + 12n\sigma\tau + 6lp^2\sigma + 6\,n\rho\tau^2 + 4\,j\sigma\tau^3,
\end{align*}
which, putting $0^4$ for $f$, may also be written
\[
\Delta\,.\,0^4 = (4k,\;4j,\;12l,\;6g,\;4l,\;12m,\;12n,\;4j,\;6l,\;6n,\;4j\!\,\Bigl)
\]
\[
(\rho^4,\,\rho^3\sigma,\,\rho^3\tau,\,\rho^2\sigma^2,\,\rho^2\tau^2,\,\rho\sigma^2\tau,\,\rho\sigma\tau,\,\rho\sigma\tau^2,\,\rho^2\tau,\,\sigma^4,\,\sigma^3\tau,\,\sigma^2\tau^2,\,\sigma\tau^3),
\]
and I join to it the series of equations
\begin{align*}
\Delta\,.\,0^3\,1 &= (0,\;3k,\;3j,\;3m,\;3l,\;3lm,\;3lm,\;3m,\;3lm,\;3l,\;l\!\,\Bigl)\,(\rho^3,\,\rho^2\sigma,\,\rho^2\tau,\,\rho\sigma^2,\,\rho\sigma\tau,\,\rho\tau^2,\,\sigma^3,\,\sigma^2\tau,\,\sigma\tau^2,\,\tau^3),\\
,,\;0^3\,2 &= (k,\;3l,\;3m,\;6m,\;3m,\;3l,\;3l\!\,\Bigl)\, ,\\
,,\;0^3\,3 &= (j,\;3l,\;3g,\;3p,\;3p,\;3l\!\,\Bigl)\, ,
\end{align*}

% ========================================================================
% PDF p. 157 = book p. 137
% ========================================================================
\medskip\noindent\textit{[book p.~137]}\quad
\begin{align*}
\Delta\,.\,0^2\,1^2 &= (0,\;2k,\;2k,\;2l,\;q\!\,\xi\,p,\,p,\,p,\,p),\quad \text{,, ,, }\\
,,\;0^2\,1\,2 &= (k,\;2k,\;3l,\;3lm,\;n\!\,\xi\;\;,\;,\;,\;,,\;,\;,)\,,\\
,,\;0^2\,1\,3 &= (j,\;2l,\;2k,\;l,\;k\!\,\xi\,p,\,p,\,p,\,p),\quad \text{,, ,, }\\
,,\;0^2\,2^2 &= (k,\;2m,\;2n,\;l,\;p\!\,\xi\;\;,\;,\;,\;,,\;,\;,)\,,\\
,,\;0^2\,2\,3 &= (l,\;2m,\;2m,\;f,\;l\!\,\xi\,p,\,p,\,p,\,p),\quad \text{,, ,, }\\
,,\;0^2\,3^2 &= (q,\;2m,\;2p,\;p,\;2l,\;0\!\,\xi\;\;,\;,\;,\;,,\;,\;,)\,,\\
\Delta\,.\,0\,1^3 &= (0,\;h,\;l,\;l\!\,\xi\,\rho,\,\sigma,\,\tau),\\
,,\;0\,1^2\,2 &= (h,\;l,\;l\!\,\xi\;\,\ldots\,),\\
,,\;0\,1\,2^2 &= (l,\;m,\;n\!\,\xi\;\,\ldots\,),\\
,,\;0\,1^2\,3 &= (j,\;l,\;m\!\,\xi\;\,\ldots\,),\\
,,\;0\,1\,2\,3 &= (l,\;m,\;n\!\,\xi\;\,\ldots\,),\\
,,\;0\,1\,3^2 &= (q,\;n,\;p\!\,\xi\;\,\ldots\,),\\
,,\;0\,2^3 &= (k,\;l,\;l\!\,\xi\;\,\ldots\,),\\
,,\;0\,2^2\,3 &= (l,\;m,\;n\!\,\xi\;\,\ldots\,),\\
,,\;0\,2\,3^2 &= (m,\;f,\;l\!\,\xi\;\,\ldots\,),\\
,,\;0\,3^3 &= (q,\;k,\;0\!\,\xi\;\,\ldots\,),
\end{align*}
which will be convenient in the sequel.

\medskip\noindent\textit{Major Function---The Fixed Curve a Cubic.}\quad Art.\ No.\ 53.

\medskip\noindent\textbf{53.}\quad It has been already seen that a simple particular form is $\tfrac{1}{3}\Omega = \bdet{012}$: and that the general form is $\Omega = \Omega' + K\,.\,012$.

\medskip\noindent\textit{Major Function---The Fixed Curve a Quartic.}\quad Art.\ No.\ 54.

\medskip\noindent\textbf{54.}\quad It is to be shown that a particular form is
\[
\tfrac{1}{2}\Omega' = \frac{-\,0\,1^3\,.\,02^2 + 01^{2}\,2\,.\,012 + 01\,2\,.\,1^{2}\,2}{1^{2}\,2}.
\]
In fact, by the foregoing values of $\Delta\,.\,01^3$, \&c., the numerator of this expression, multiplied by $\Delta!$, is $=$
\begin{align*}
&\quad -(ha + jr)(la + fr)\\
&\quad + (kp + ra + lr)(c\rho + l\sigma + nr)\\
&\quad + r(h\rho^{2} + 2r\rho\sigma + 2lr\sigma + lr\sigma + lnr\sigma + n\sigma^{3}),
\end{align*}
which is
\[
= 2lr\rho^{2} + 3r^{2}\rho\sigma + (lm - jk + 3lr)\rho\sigma + 2lr\sigma^{2} + r(-fh + kl + 3n\sigma\tau + r\tau^{2},
\]
and this, for $\tau = 0$, becomes
\[
= r\,(2k\rho^{2} + 3r\rho\sigma + 2l\sigma^{2}).
\]
C.\ XII. \hfill 19

% ========================================================================
% PDF p. 158 = book p. 138
% ========================================================================
\medskip\noindent\textit{[book p.~138]}\quad Hence for $\tau = 0$, we have
\[
\Omega' = \frac{1}{\Delta!}\,(4k\rho^{2} + 6r\rho\sigma + 4l\sigma^{2}),
\]
that is,
\[
= \frac{1}{\Delta}\,[4\,.\,1^{3}\,2\,.\,\rho^{2} + 6\,.\,1^{2}\,2^{2}\,.\,\rho\sigma + 4\,.\,1\,2^{3}\,.\,\sigma^{2}];
\]
and $\Omega'$ is thus a form of the major function $(x,\,y,\,z)_{2}^{2}$. Of course the general form is $\Omega = \Omega' + (x,\,y,\,z)\,.\,012$.

\medskip\noindent\textit{Syzygy of the Major Function.}\quad Art.\ No.\ 55.

\medskip\noindent\textbf{55.}\quad Writing now $(x,\,y,\,z)_{k}^{n-3} = \Omega_{ik}$; and taking on the fixed curve a new point $3$, consider the like functions $\Omega_{23}$ and $\Omega_{31}$: it is to be shown that we have identically
\[
\Omega_{23}\,.\,011 + \Omega_{31}\,.\,012\,.\,023 + \Omega_{12}\,.\,031 - (123)^{n}\,.\,011\,.\,012\,(x,\,y,\,z)^{n-4},
\]
where $(x,\,y,\,z)^{n-4} = 0$ is a properly determined minor function; or, considering herein $0$ as a point on the fixed curve and writing therefore $f = 0$, the equation is
\[
\frac{\Omega_{23}}{023}\,.\,011 + \frac{\Omega_{31}}{031}\,.\,012 + \frac{\Omega_{12}}{012}\,.\,(x,\,y,\,z)^{n-4}.
\]

\medskip\noindent\textbf{56.}\quad Write for a moment $X = \Omega_{23}\,.\,011\,.\,012 + \Omega_{31}\,.\,023\,.\,031 - (123)^{n}\,(x,\,y,\,z)^{n-4}$, then $k$ being an arbitrary coefficient, we have $X - kf = 0$, a curve of the order $n$, passing through the points $1$, $2$, $3$, and the residues of $1$, $2$; in fact, take for any other point in the curve $\Omega_{12}\,.\,012 = 0$, $033 = 0$, and therefore $X = 0$ a point on the curve. Again, at any residue of $2$, $3$, we have $\Omega_{23} = 0$, $023 = 0$, and therefore $X = kf = 0$ will have a dp at $1$; and clearly by symmetry, it will also have a dp at $2$, and at $3$; the theorem is thus proved.

\medskip\noindent\textbf{57.}\quad Taking an arbitrary point $a$ coordinates $(x_a,\,y_a,\,z_a)$, and writing
\[
D = x_a\,\frac{d}{dx} + y_a\,\frac{d}{dy} + z_a\,\frac{d}{dz},
\]
\bigskip
\hrule
\smallskip
\noindent\textsuperscript{*}\,This is the differential theorem corresponding to $C$, and $G$.'s integral theorem, p.\ 30, viz.\ this is $\Omega_{23} + \Omega_{31} + \Omega_{12} = 1$, a sum of three integrals of the third kind=an integral of the first kind.

% ========================================================================
% PDF p. 159 = book p. 139
% ========================================================================
\medskip\noindent\textit{[book p.~139]}\quad we have to find $k$, so that the curve $D(X - kf) = 0$ shall pass through the point $1$. Observing that $D023 = a23$, \&c., we have
\begin{align*}
D(X - kf) &= D\Omega_{23}\,.\,031\,.\,012 + \Omega_{23}\,(031\,.\,a12 + a31\,.\,012)\\
&\quad + a23\,(\Omega_{31}\,.\,012 + \Omega_{12}\,.\,031)\\
&\quad + 023\,[\Omega_{31}\,.\,a12 + \Omega_{31}\,.\,a31 + D\Omega_{31}\,.\,012 + D\Omega_{12}\,.\,031]\\
&\quad - kDf,
\end{align*}
and, to make the curve pass through $1$, writing herein $0 = 1$, we have
\[
0 = 123\,(\Omega_{12}^{2}\,.\,a12 + \Omega_{1,a}\,.\,a31) - k\,(Df)_{1},
\]
where the suffixe (1) denotes that we are in $\Omega_{23}$, \&c., and $f$ respectively to write $0 = 1$. We have here $\Omega_{1,a} = a\,.\,1^{n-3}\,3$, $\Omega_{12,1} = n\,.\,1^{n-3}\,2$, $(Df)_{1} = a\,.\,1^{n-1}$, and the equation thus is
\[
n\,.\,123\,(1^{n-2}\,3\,.\,a12 + 1^{n-2}\,2\,.\,a31) = ka\,.\,1^{n-1}.
\]

But we have identically $1^{n-2}\,1\,.\,a23 + 1^{n-2}\,2\,.\,a31 + 1^{n-2}\,3\,.\,a12 = 1^{n-1}\,.\,123$, where $1^{n-1} = 1^{n}$ is, in fact, $= 0$; the factor $1^{n-2}\,a$ thus divides out, and the equation becomes $k = (123)^{n}$, viz.\ $k$ having this value, the curve $X - kf = 0$ will have a dp at $1$; and clearly by symmetry, it will also have a dp at $2$, and at $3$; the theorem is thus proved.

\medskip\noindent\textit{The Syzygy, Fixed Curve a Cubic.}\quad Art.\ No.\ 58.

\medskip\noindent\textbf{58.}\quad The syzygy may be verified independently in the case where the fixed curve is a cubic. Observe that the syzygy, if satisfied for any particular form of $\Omega$, will be generally satisfied; we may therefore take $\tfrac{1}{3}\Omega_{12} = \bdet{012}$. Writing thus
\[
\frac{\tfrac{1}{3}\Omega_{23}}{012}\,.\,\frac{\bdet{012}}{012} = [012]\;\text{suppose,}
\]
and taking $0$ to be a point on the cubic curve, we ought to have $[023] + [031] + [012] = a$ constant; the value of this constant comes out to be $= [123]$, and the syzygy in its complete form thus is
\[
[023] + [031] + [012] = [123].
\]
We have
\[
\Delta\,\bdet{023},\;\Delta\,\bdet{031},\;\Delta\,\bdet{012} = lp + f\sigma + i\tau,\;jp + h\sigma + g\tau,\;h\rho + l\sigma + l\tau,
\]
and the equation thus is
\[
\frac{lp + f\sigma + i\tau}{\rho} + \frac{jp + h\sigma + g\tau}{\sigma} + \frac{h\rho + l\sigma + l\tau}{\tau} - l = 0;
\]
this, multiplied by $\rho\sigma\tau$, becomes
\[
h\rho^{2}\sigma + jp^{2}\tau + 3l\rho\sigma\tau + g\sigma^{2}\tau + f\sigma\tau^{2} + i\sigma\tau^{2} = 0,
\]
which is, in fact, $\tfrac{1}{3}f = 0$, the equation of the cubic curve.

Observe that the new symbol $[012]$ is, in virtue of its determinant denominator, an alternate function, $[012] = -\,[102]$, $[012] = [201] = [201]$. The syzygy is a relation between any four points $1$, $2$, $3$, $0$ of the curve, and it may be also expressed in the form
\[
[123] - [230] + [301] - [012] = 0.
\]
\hfill $18$--$2$

% ========================================================================
% PDF p. 160 = book p. 140
% ========================================================================
\medskip\noindent\textit{[book p.~140]}\quad \textit{The Syzygy, Fixed Curve a Quartic.}\quad Art.\ No.\ 59.

\medskip\noindent\textbf{59.}\quad Taking $\Omega_{12}$ as before, we have
\[
\tfrac{1}{2}\Omega_{12} = \frac{-\,01^{3}\,.\,02^{2} + 01^{2}\,2\,.\,012 + 01\,2\,.\,1^{2}\,2}{012\,.\,1^{2}\,2} = [012]\;\text{suppose;}
\]
and taking $0$ to be a point on the quartic curve, we ought to have
\[
[023] + [031] + [012] = (x,\,y,\,z)^{0},\;\text{a linear function of }(x,\,y,\,z),
\]
or, what is the same thing, considering the left-hand side as expressed in terms of $\rho$, $\sigma$, $\tau$, the sum should be
\[
= (\rho,\,\sigma,\,\tau)^{0},\;\text{a linear function of }(\rho,\,\sigma,\,\tau).
\]

By a preceding formula we have
\begin{align*}
[0\,1\,2] &= \frac{1}{\Delta^{2}\rho\sigma}\,\bigl[2lr\rho^{2} + 3r^{2}\rho\sigma + km - jk + 3lr)\rho\sigma\\
&\quad + 2lr\sigma^{2} + (-\,fh + kl + 3n\sigma\tau)\tau^{2}\bigr],
\end{align*}
which is
\[
= \frac{1}{\Delta!}\left\{(2k + \frac{km - jk}{r})\,\rho + \left(3m + \frac{-fh + kl}{r}\right)\sigma + \left(n + \frac{-fj + lm}{r}\right)\tau\right\} + \frac{1}{\Delta!}\,\frac{2k\rho^{2} + 3r\rho\sigma + 2l\sigma^{2}}{\tau}.
\]

And hence, forming the sum $[023] + [031] + [012]$, we have first a fractional part which is found to be integral, viz.\ this is
\[
\frac{1}{\Delta^{2}}\left\{\frac{2fr^{2} + 3r\rho\sigma + 2ir^{2}}{p} + \frac{2gp^{2} + 3r\rho\sigma + 2lp^{2}}{q} + \frac{2kp^{2} + 3r\rho\sigma + 2l\sigma^{2}}{r}\right\},
\]
\[
= \frac{1}{\Delta^{2}\rho\sigma\tau}\left\{2hp\sigma + 2lp\rho + 3l\rho^{2}\sigma + 3gp^{2}\sigma + 2lp\sigma\tau + 2k\rho\sigma\tau + 2k\rho\tau\right\},
\]
\[
= \frac{1}{\Delta!\rho\sigma\tau}\bigl[\tfrac{1}{2}\Delta!\,f - 6lp\sigma\tau - 6n\sigma\tau - 6n\sigma\tau\bigr],
\]
or since $f = 0$, this is
\[
= \frac{1}{\Delta}\,(-\,6l\rho - 6n\sigma - 6n\tau).
\]

We then have integral terms which are at once deduced from the above integral terms of $0^{1}\,2$; collecting the several terms, we find
\[
[023] + [031] + [012] =
\]
\[
\frac{1}{\Delta}\left\{\rho\,\left(l + \frac{mn - gk}{p} + \frac{in - gh}{q} + \frac{kn - gh}{r}\right) + \sigma\,\left(m + \frac{fn - hi}{p} + \frac{lh - lk}{q} + \frac{kl - fh}{r}\right) + \tau\,\left(n + \frac{lm - fj}{p} + \frac{gl - ij}{q} + \frac{lm - fj}{r}\right)\right\},
\]
which is the required result.

% ========================================================================
% PDF p. 161 = book p. 141
% ========================================================================
\medskip\noindent\textit{[book p.~141]}\quad \textit{Preparation for the Conversion---The Symbol $\partial$.}\quad Art.\ Nos.\ 60 to 63.

\medskip\noindent\textbf{60.}\quad I use $\partial$ as the symbol of a quasi-differentiation, viz.\ $U$ being any function of $(x,\,y,\,z)$, $\partial U$ denotes $\dfrac{1}{d\omega}$ multiplied by the differential $\dfrac{dU}{dx}\,dx + \dfrac{dU}{dy}\,dy + \dfrac{dU}{dz}\,dz$; in such a differential, the increments $dx$, $dy$, $dz$ do not in general present themselves in the combinations $ydz - zdy$, $zdx - xdz$, $xdy - ydx$; but they will do so if $U$ is a function of the degree zero in the coordinates $x$, $y$, $z$, that is, if $U$ be the quotient of two homogeneous functions of the same degree(s); and this being so, we say by the equations
\[
\frac{ydz - zdy}{\dfrac{df}{dx}} = \frac{zdx - xdz}{\dfrac{df}{dy}} = \frac{xdy - ydx}{\dfrac{df}{dz}}\;= d\omega
\]
get rid of the increments, and $\partial U$ will denote a function of $(x,\,y,\,z)$ derived in a definite manner from the function $U$. The symbol $\partial$ will be used only in the case in question of a function of the degree zero. Of course $\partial_1$ will denote the like operation in regard to $(x_1,\,y_1,\,z_1)$, and so $\partial_2$, \&c.; and we may for greater clearness write $\partial_0$ in place of $\partial$.

\medskip\noindent\textbf{61.}\quad Consider then $\partial \dfrac{P}{Q}$, where $P$, $Q$ are functions $(x,\,y,\,z)^m$ of the same degree; we have
\[
\partial\,\frac{P}{Q} = \frac{1}{Q^2 d\omega}\,(QdP - PdQ),
\]
and then
\[
dP = \frac{dP}{dx}\,dx + \frac{dP}{dy}\,dy + \frac{dP}{dz}\,dz,\quad \frac{1}{m}P = \frac{dP}{dx}\,x + \frac{dP}{dy}\,y + \frac{dP}{dz}\,z,
\]
with the like formula for $Q$. Substituting, we find
\[
\partial\,\frac{P}{Q} = \frac{1}{mQ^2 d\omega}\left\{\frac{d(Q,\,P)}{d(y,\,z)}\,(ydz - zdy) + \frac{d(Q,\,P)}{d(z,\,x)}\,(zdx - xdz) + \frac{d(Q,\,P)}{d(x,\,y)}\,(xdy - ydx)\right\},
\]
that is,
\[
\partial\,\frac{P}{Q} = \frac{1}{mQ^2}\,\left\{\frac{df}{dx}\,\frac{d(Q,\,P)}{d(y,\,z)} + \&\text{c.}\right\} = \frac{1}{mQ^2}\,\frac{d(f,\,Q,\,P)}{d(x,\,y,\,z)},\;=\;\frac{1}{mQ^2}\,J\,(f,\,Q,\,P),
\]
or say
\[
\partial\,\frac{P}{Q} = -\,\frac{1}{mQ^2}\,J\,(P,\,Q,\,f).
\]

\medskip\noindent\textbf{62.}\quad As an example, consider
\[
\partial\,[012],\;= -\,\frac{1}{(012)^2}\,J\,(\bdet{012},\;012,\;f).
\]
The determinant is
\[
\left|\begin{array}{ccc}
\dfrac{d}{dx}\,\bdet{012}, & y_1z_z - y_2z_1, & \dfrac{df}{dx}\\[4pt]
\dfrac{d}{dy}\,\bdet{012}, & z_1z_2 - z_2x_1, & \dfrac{df}{dy}\\[4pt]
\dfrac{d}{dz}\,\bdet{012}, & x_1y_2 - x_2y_1, & \dfrac{df}{dz}
\end{array}\right|
\]

% ========================================================================
% PDF p. 162 = book p. 142
% ========================================================================
\medskip\noindent\textit{[book p.~142]}\quad and the coefficient herein of $\dfrac{d}{dz}\,\bdet{012}$ is $(x_1z_2 - x_2z_1)\dfrac{df}{dy} - (x_1y_2 - x_2y_1)\dfrac{df}{dx}$, which is
\[
= x_1\left(x_2\,\frac{df}{dx} + y_2\,\frac{df}{dy} + z_2\,\frac{df}{dz}\right) - x_2\left(x_1\,\frac{df}{dx} + y_1\,\frac{df}{dy} + z_1\,\frac{df}{dz}\right),\;= 3\,(0^{2}\,.\,x_1 - 0^{2}\,.\,x_2),
\]
and so for the other terms.

The determinant is thus
\[
= 3\,\left[0^{2}\,1\,\left(x_1\,\frac{d}{dx} + y_1\,\frac{d}{dy} + z_1\,\frac{d}{dz}\right)\bdet{012} - 0^{2}\,2\left(x_2\,\frac{d}{dx} + y_2\,\frac{d}{dy} + z_2\,\frac{d}{dz}\right)\bdet{012}\right]
\]
say this is
\[
= 3\,[0^{2}\,1\,.\,\bdet{D} - 0^{2}\,2\,\bdet{D}]\,\bdet{012}.
\]

But we have $\bdet{D}\,\bdet{012} = 1^{2}\,2$, $\bdet{D}\,\bdet{012} = 1\,2^{2}$, and the determinant is then $= 3\,(0^{2}\,.\,1^{2}\,2 - 0^{2}\,.\,1\,2^{2})$, whence finally, writing $\partial_0$ instead of $\partial$,
\[
\partial_0\,[012] = -\,3\,.\,\frac{0^{2}\,.\,1^{2}\,2 - 0^{2}\,.\,1\,2^{2}}{(012)^{2}}.
\]

\medskip\noindent\textbf{63.}\quad By cyclical interchange of the $0$, $1$, $2$, we have
\[
\partial_1\,[012] = -\,3\,.\,\frac{1^{2}\,.\,0^{2}\,2 - 0^{2}\,.\,02^{2}}{(012)^{2}},
\]
\[
\partial_2\,[012] = -\,3\,.\,\frac{02^{2}\,.\,01^{2} - 1^{2}\,.\,0^{2}\,1}{(012)^{2}},
\]
and thence adding, we find
\[
(\partial_0 + \partial_1 + \partial_2)\,[012] = 0,
\]
an important property, which, joined to the equation before obtained,
\[
[023] + [031] + [012] = [123],
\]
completes the theory of the function $[012]$.

\medskip\noindent\textit{Conversion of the Major Function (Interchange of Limits and Parametric Points).}\quad Art.\ No.\ 64.

\medskip\noindent\textbf{64.}\quad Write in general
\[
\frac{(x,\,y,\,z)^{n-3}_{k}}{012} = Q_{k,\,01}.
\]
$Q_{k,\,01}$ is an alternate function in regard to the points $1$, $2$ ($Q_{k,\,01} = -\,Q_{k,\,10}$), and in regard to the coordinates of the points $0$, $1$, $2$, it is rational, but not integral, of the degrees $n - 3$, $0$, $0$ respectively; it can therefore be operated upon with $\partial_0$ or $\partial_1$, but (except in the case $n = 3$) not with $\partial_2$.

% ========================================================================
% PDF p. 163 = book p. 143
% ========================================================================
\medskip\noindent\textit{[book p.~143]}\quad The conversion relates not to the general major function $(x,\,y,\,z)^{n-3}_{k}$, but to this function with the arbitrary constants properly determined, and consists in a relation between two functions $Q_{k,\,01}$ and $Q_{k,\,1k}$ (each of them a function of three out of four arbitrary points $1$, $2$, $4$, $5$ on the fixed curve), viz.\ the conversion is
\[
\partial_2\,Q_{k,\,01} = \partial_4\,Q_{k,\,4k};
\]
an equation which may be written in four different forms, viz.\ we may in the form written down interchange $1$, $2$ and also $4$, $5$.\textsuperscript{*}

The determination of the constants is a very peculiar one, inasmuch as it is not algebraical; viz.\ in the case of the cubic curve, about to be considered, it appears that $Q_{k,\,01}$ contains the term $\int_2^1 d\omega\,\bigl\{0\bdet{36}\bigr\}$, which is a transcendental function of the coordinates of the parametric points $1$ and $2$.

\medskip\noindent\textit{The Conversion, Fixed Curve a Cubic.}\quad Art.\ No.\ 65.

\medskip\noindent\textbf{65.}\quad We may write $Q_{k,\,01} = [012] + K$, where $K$ is a constant, that is, it is independent of the point $0$, but depends on the parametric points $1$ and $2$. I assume $K$ to be properly determined, and give an \textit{a posteriori} verification of the equation
\[
\partial_2\,Q_{k,\,01} = \partial_4\,Q_{k,\,4k}.
\]
The value is $K = \int_2^1 d\omega\,\bigl\{0\bdet{36}\bigr\} - [123]$, where $3$, $4$ are arbitrary points on the cubic curve, and where in the definite integral, regarded as an integral $\int U\,du$ with a current variable $u$, the meaning is that this variable has as the limits the values $u_1$, $u_2$, which belong to the points $1$ and $2$ respectively: a fuller explanation might be proper, but the investigation will presently be given in a form not depending on any integral at all.

Substituting for $K$ its value, we have
\[
Q_{k,\,01} = [012] + \int_2^1 d\omega\,\bigl\{0\bdet{36}\bigr\} - [123],
\]
or, as this may also be written,
\[
= -\,[023] - [031] + \int_2^1 d\omega\,\bigl\{0\bdet{36}\bigr\}.
\]
We have thence
\[
\partial_2 Q_{k,\,01} = -\,\partial_2\,[031] + \int_2^1 d\omega\,\bigl\{0\bdet{36}\bigr\},
\]
\bigskip
\hrule
\smallskip
\noindent\textsuperscript{*}\,The meaning of the property is better seen from the integral form: $Q_{k,\,01}$ is a function of the points $0$, $1$, $2$ and $Q_{k,\,4k}$ the like function of the points $0$, $1$, $4$, $5$, and these in their meaning (alluded to in the heading) that there exists for the integral of the third kind $\Omega$ a canonical form $\Pi$, such that $\int_0^{(\alpha\beta)} d\omega\, \bigl\{0\bdet{36}\bigr\}_{\alpha\beta} = \int_a^{(\alpha\beta)} d\omega\, Q_{k,\,a\beta}$ which equation operated upon with $\partial_0$ gives the formula of the text. And there is then the existing (alluded to in the heading) that there exists for the integral of the third kind a canonical form $\Pi$, in which the symmetry, expressed by the interchange of the limits and the parametric points. The expression for $Q_{k,\,a\beta}$ mentioned further on in the text for the case, must cover a cubic, shows that in this case the canonical form of the integral of the third kind is $\int_a^b d\omega\,\bigl\{[026] + \left(\int_a^b d\omega\,\{[036] - [126]\}\right)\bigr\}$.

% ========================================================================
% PDF p. 164 = book p. 144
% ========================================================================
\medskip\noindent\textit{[book p.~144]}\quad and consequently
\[
\partial_2 Q_{k,\,01} = -\,\partial_2\,[431] + \partial_2\,[136],
\]
\[
\partial_4 Q_{k,\,4k} = -\,\partial_4\,[134] + \partial_4\,[436];
\]
hence, observing that $[431] = -\,[134]$, \&c., we have
\[
\partial_2 Q_{k,\,01} - \partial_4 Q_{k,\,4k} = -\,\partial_2\,[134] + \partial_2\,[136] - \partial_4\,[436],
\]
\[
= -\,\partial_2\,[134] + \partial_2\,[136] - \partial_4\,[436],
\]
which, observing that we have $\partial_2\,[641] = 0$, is
\[
= \partial_2\,([136] - [364] + [641] - [413]),\;= 0,
\]
the required theorem.

To avoid, in the proof, the use of the integral sign, we have only to consider the required function $Q_{k,\,01}$ as given by the foregoing differential formula
\[
\partial_2 Q_{k,\,01} = -\,\partial_2\,[031] + \partial_2\,[136],
\]
for we have then the values of $\partial_1 Q_{k,\,01}$ and $\partial_2 Q_{k,\,01}$: the rest of the proof the same as before.

\medskip\noindent\textit{The Conversion, Fixed Curve a Quartic.}\quad Art.\ Nos.\ 66 to 73.

\medskip\noindent\textbf{66.}\quad We have
\[
Q_{k,\,01} = [012] + (x,\,y,\,z),
\]
where $(x,\,y,\,z)$ is a linear function of $(x,\,y,\,z)$, but depending also on the parametric points $1$ and $2$, which has to be determined so as to satisfy the conversion equation
\[
\partial_2 Q_{k,\,01} = \partial_4 Q_{k,\,4k}.
\]

Observing that we have $[023] + [031] + [012] = $ a linear function of $(x,\,y,\,z)$, the linear function $(x,\,y,\,z)$ of $Q_{k,\,01}$ may be taken to be $= \Theta_{k,\,12} = [023] - [031] - [012]$; that is, we may assume
\[
Q_{k,\,01} = [012] + \Theta_{k,\,12} = (-[023] + [091]) + [012],
\]
\[
= -\,[023] - [031] + \Theta_{k,\,12},
\]
where $\Theta_{k,\,12}$ is a linear function of $(x,\,y,\,z)$, but depending also on the points $1$ and $2$, which has to be determined. We have
\[
Q_{k,\,01} = [012] + \Theta_{k,\,12} + \partial_1\Theta_{k,\,12},
\]
and thence
\[
\partial_1 Q_{k,\,01} = -\,\partial_1\,[431] + \partial_1\Theta_{k,\,12},
\]
\[
\partial_2 Q_{k,\,01} = -\,\partial_2\,[134] + \partial_2\Theta_{k,\,12},
\]
giving an equation for $\Theta$,
\[
\partial_1\Theta_{k,\,12} - \partial_2\Theta_{k,\,12} = \partial_2\,[431] - \partial_2\,[134];
\]
here $4$ is an arbitrary point of the quartic, and we may instead of it write $0$, the equation thus becoming
\[
\partial_1\Theta_{k,\,12} - \partial_2\Theta_{k,\,12} = \partial_2\,[031] - \partial_2\,[130].
\]

% ========================================================================
% PDF p. 165 = book p. 145
% ========================================================================
\medskip\noindent\textit{[book p.~145]}\quad \textbf{67.}\quad Of the terms on the left-hand side, the first is a linear function of $(x,\,y,\,z)$, or say it is an integral function of $0$, and the second is a linear function of $(x_1,\,y_1,\,z_1)$, or say it is an integral function $1'$; the first depends on the parametric points $1$, $3$, and $\partial_1(0,\,1,\,3)$ is a known function, suppose, as regards the coordinates $(x,\,y,\,z)$ of the point $0$, replacing the corresponding point $3$, $4$ or $5$ by the new arbitrary point $6$, $7$ or $8$. Further $045$, \&c., denote determinants as above; so that in $H$ each of the last three terms is, in fact, as regards the point $0$, a mere linear function of the coordinates $(x,\,y,\,z)$ of this point.

We have $\partial_1(1\,;\,3,\,2;\,6,\,4,\,5) = 13'\bdet{246}$, and hence this function multiplied by $\dfrac{045}{345}$ is $= 01'\bdet{2456}$; and so for the third and fourth terms of $H$: thus each of the four terms of $H$ is $= 01'\bdet{2456789}$, of the degree $1$ in the coordinates of the points of the other points $2$, $3$, $4$, $5$, $6$, $7$, $8$ respectively, but of the degree $0$ in regard to the coordinates of the point $1$. The before-mentioned function $\Omega\,(0;\;1,\;2;\;3,\;4,\;5)$ is a function satisfying these conditions $1\Omega$ as $= s\,.\,012\,1^{2}\,2 +$ arbitrary linear function $(x,\,y,\,z)$, of $(x_1,\,y_1,\,z_1)$, of the points of of $(x,\,y,\,z)$, which presents itself in the case of the cubic curve.

\medskip\noindent\textit{Nöther's conversion-theorem consists herein, that the function}
\[
H\,(0;\;1,\;2;\;3,\;4,\;5;\;6,\;7,\;8)
\]
is unaltered by the interchange of the two points $0$, $1$; or putting for shortness
\[
H\,(0;\;1,\;2;\;3,\;4,\;5;\;6,\;7,\;8) = H_1\,(0),
\]
the theorem is
\[
H_1\,(0) = H_0\,(1).
\]

\medskip\noindent\textbf{74.}\quad We have, No.~59,
\[
\tfrac{1}{2}\Omega_{12} = \frac{-\,01^{3}\,.\,02^{2} + 01^{2}\,2\,.\,012 + 01\,2\,.\,1^{2}\,2}{012\,.\,1^{2}\,2},\;= [012];
\]
or as for greater simplicity I write it
\[
= 012,
\]
viz.\ $012$ is now written instead of $[012]$ to denote the function just given as the value of $\dfrac{\tfrac{1}{2}\Omega_{12}}{012}$. $\Omega_{12}$ is thus $= r\,.\,012\,.\,012$, viz.\ this is a particular form of $\Omega_{12}$ satisfying the conditions that $\Omega_{12} = 0$ is a conic passing through the residues of the points $1$, $2$, and such that $\Omega_{12}$ on writing therein $(x_1,\,y_1,\,z_1)$ for $(x,\,y,\,z)$ becomes $= s\,.\,012\,1^{2}\,2 +$ arbitrary linear function $(x,\,y,\,z)$. The before-mentioned function $\Omega\,(0;\;1,\;2;\;3,\;4,\;5)$ is a function satisfying these conditions $1\Omega$ as $= s\,.\,012\,1^{2}\,2 +$ arbitrary linear function $(x,\,y,\,z)$, of $(x_1,\,y_1,\,z_1)$, of the points of of $(x,\,y,\,z)$, which presents itself in the case of the cubic curve.

% ========================================================================
% PDF p. 166 = book p. 146
% ========================================================================
\medskip\noindent\textit{[book p.~146]}\quad we find, by a process such as that of No.\ 62,
\[
\partial_1\,[0\,1\,3] = \frac{1}{a}\,\left\{-\,0\,1^{3}\,.\,01^{2}\,.\,02 + 09^{2}\,.\,j\right\}
\]
\[
+ \frac{1}{q}\,\bigl\{2\,.\,01^{3}\,(03^{2}\,.\,01^{2}\,2 + (013)^{2}\bigr\}_{\beta}
\]
\[
+ \frac{1}{q}\,\bigl\{(2\,.\,01^{3}\,.\,1\,2)\,.\,(01\,2\,.\,03^{2}\,+\,2\,(2\,.\,013\,.\,01^{2}\,.\,01^{2}\,2\,+\,01^{2}\,.\,03)\bigr\}_{\beta}
\]
\[
+ \frac{1}{q^{2}}\,\bigl\{-\,2\,.\,01^{3}\,(-\,01^{2}\,.\,01^{2}\,+\,01^{2}\,.\,03^{2})\bigr\}\,.
\]

Substituting herein the value $01^{3} = \dfrac{1}{a^{3}}(ha + jr)$, \&c., we have $\dfrac{1}{a^{3}}$ multiplied by a cubic function $(\rho,\,\sigma,\,\tau)$; writing down very simply the integral terms, and then the others, we have
\[
\partial_1\,[013] = \frac{1}{a^{3}}\bigl\{(-\,6hn + 3jm) + \frac{1}{q}\,(8hkl - 3qjp + 3jp + 2jh)\rho\sigma
\]
\[
+ \frac{1}{q^{2}}\,\bigl(\tfrac{15}{4}h^{2} - 7hp + 8jm\bigr)\rho^{2} + \frac{1}{q}\,(-\,2pjk + 8jh - 5jp + 4h^{2})\rho\sigma + \frac{1}{q^{2}}\,(2j^{2}hl - 2pj n - 2pjm + 2j^{2}n)\bigr\}
\]
(say this linear function of $\rho$, $\sigma$, $\tau$ is $= \square$)
\[
+ \frac{1}{\Delta^{2}p^{2}}\,\bigl[ar^{3}\,.\,2j^{2} + p^{2}\sigma\rho + p^{2}\sigma\rho + p\sigma\tau + 6jp + 6jp + 6jn\sigma)\bigr] - *
\]

\medskip\noindent\textbf{70.}\quad The expression of $\partial_1\,[103]$ is deduced from this by the interchange of $0$, $1$: and I write
\[
\partial_1\,[013] - \partial_1\,[103] = \square - *
\]
\[
+ \frac{1}{a^{3}\Delta!p}\bigl\{p^{2}\rho^{2}\,2j^{2} + p^{2}\sigma\,3jh + p^{2}\tau\,3l + p^{2}\sigma\,3lk + p\rho\sigma + p\sigma\tau + p^{2}\rho + p\sigma\tau + 2j\sigma\tau\,3jn\bigr\}
\]
\[
- \Delta^{2}\,\Delta\,.\,2\,(09\bar{1})\sigma - \Delta!\sigma\,(-\,3\,.\,01^{2}\,.\,02^{2} + 0\bar{2}^{2}\,.\,025) - \Delta!\,3\,.\,01^{3}\,.\,01^{2}
\]
\[
+ \Delta\sigma\,(-\,2\,.\,01^{3}\,.\,02 + 0\bar{2}^{2}\,(02 + 01\bar{2})) - 2\,.\,01^{3}\,.\,01^{2}
\]
where, and in what follows, the $*$ denotes the function included immediately to the left of it, interchanging therein the $0$, $1$. It will be observed that the $\square$, and linear function of $(\rho,\,\sigma,\,\tau)$, that is, of $(x,\,y,\,z)$, is a series of the required function $\partial_1\,(0;\;1,\;3)$: the remaining portion has to be denoted by means of the expressions for $\Delta^{2}(013)$, \&c., in terms of $\rho$, $\sigma$, $\tau$.

\medskip\noindent\textbf{71.}\quad We obtain
\[
\partial_1\,[013] - \partial_1\,[103] = \square - *
\]
\[
+ \frac{1}{\Delta!}\,\bigl[\sigma\,(2fj - 3kp - 3hq + 15lm - 9nr) + (-\,3gr - 3jp + 9mq)\bigr]
\]
\[
+ \frac{1}{\Delta!p}\,\bigl[\sigma^{2}\,(12fk - 12hn + 18m^{2} - 9pr) + \sigma\tau\,(6fp - 6jr + 18nq) + \tau^{2}\,.\,4n\sigma\bigr]
\]
\[
+ \frac{1}{\Delta!p^{2}}\,\bigl[\rho^{3}\,(18gn - 9pkp) + \rho^{2}\sigma\,(12hp - 8tk + 9mqp) + \sigma\tau\,(4fp + 6lm)\bigr]
\]
\[
+ \frac{1}{\Delta!p^{2}}\,\bigl[\rho^{3}\,.\,4f^{2} + \rho^{2}\sigma\,.\,6fp + \rho\sigma\,.\,4f j\bigr].
\]

% ========================================================================
% PDF p. 167 = book p. 147
% ========================================================================
\medskip\noindent\textit{[book p.~147]}\quad The terms of the second line may be transformed as follows:
\[
\frac{\sigma}{\Delta!}\,(2fj - 3kp - 3hq + 18lm - 9nr)
\]
\[
= \tfrac{1}{2}\,\frac{\sigma}{\Delta!}\,(2fj - 3kp - 3hq + 18lm - 9nr) - *
\]
\[
+ \frac{1}{\Delta!p}\,\bigl[\sigma^{2}\,(-\,12fk + 12hn - 18m^{2} + 9pr) + \sigma\tau\,\bigl\{\tfrac{1}{2}fp + 3pk + \tfrac{3}{4}jr - \tfrac{3}{4}(p - 9nm)\bigr\}\bigr]
\]
\[
+ \frac{1}{\Delta!p^{2}}\,\bigl[\sigma^{3}\,(-\,20fm + 15hp) + \sigma^{2}\tau\,(-\,\tfrac{3}{2}fr + 12lk - 18np) + \sigma\tau^{2}\,.\,-\,6\sigma p\bigr]
\]
\[
+ \frac{1}{\Delta!p^{2}}\,\bigl[\sigma^{2}\tau\,(-\,10fp^{2} + a^{2}\tau - 15fp + a^{2}\tau - 4mn^{2})\,a\sigma\tau + a\sigma\tau\,a^{2}\tau + a\sigma\tau\,a\sigma\tau - 6\,p\bigr]
\]
\[
+ \frac{1}{\Delta!p^{3}}\,\bigl[\sigma^{2}\tau\,(-\,10f n^{2} + a^{2}\tau - 12 f p + 8 i k - 9 i m\rho)\,a\sigma\tau - 9 i\,p\bigr]\,,
\]
and
\[
\frac{1}{\Delta!}\,(-\,3gr - 3jp + 9mq)
\]
\[
= \tfrac{1}{2}\,\frac{1}{\Delta!}\,(-\,3gr - 3jp + 9mq) - *
\]
\[
+ \frac{1}{\Delta!p}\,\bigl[\sigma\tau\,(-\,\tfrac{3}{2}fq + \tfrac{3}{2}lr + \tfrac{3}{4}(p - 9nm) + \tfrac{1}{2}\sigma\tau\,(p - 9nm)\bigr) + \tau^{2}\,.\,-\,6\sigma p\bigr]
\]
\[
+ \frac{1}{\Delta!p^{2}}\,\bigl[\sigma^{2}\tau\,(-\,6fp + a^{2}\tau - 12fp + a^{2}\tau - 6fq + a^{2}\tau - 6fq + a^{2}\tau)\,a\sigma\tau - 6\,jp\bigr]
\]
\[
+ \frac{1}{\Delta!p^{3}}\,\bigl[\sigma^{2}\tau\,(-\,6fp + a^{2}\tau - 12fp + a^{2}\tau - 6fq + a^{2}\tau)\,a\sigma\tau - 6\,jp\bigr],
\]
substituting these values, the whole third line is destroyed, and we find
\[
\partial_1\,[013] - \partial_1\,[103] = \square - *
\]
\[
+ \tfrac{1}{2}\,.\,\frac{1}{\Delta!}\,[\sigma\,(2fj - 3hq - 3kp + 18lm - 9nr) + \tau\,(-\,3gr - 3jp + 9mq)] - *
\]
\[
+ \frac{1}{\Delta!p^{2}}\,\bigl[\sigma^{2}(-\,12fm + 6fp) + \sigma\tau(-\,9fn + 7lk - 9mp - 9np) + \sigma\tau\,a^{2}\tau\,(-\,2fp - 9np) + \sigma^{2}\tau\,.\,-\,5\sigma p\bigr]
\]
\[
+ \frac{1}{\Delta!p^{3}}\,\bigl[\sigma^{2}\,(-\,6f^{2} + \sigma^{2}\tau\,a^{2}\tau - 12fp + a^{2}\tau - 6fq + a^{2}\tau)\,a\sigma\tau - 6\sigma p\bigr]\,.
\]
Ultimately the last two lines of this expression are found to be
\[
= \frac{1}{\Delta!}\,[\rho\,(-\,2hn + 4jm + 2l^{2} - 2qr) + \sigma\,(-\,2fj + 2kp + 2hq - 10lm + 5nr) + \tau\,(2gr + hi - jp + 4hr - 6mq)] - *
\]
so that the whole is now a sum of three linear functions of $(\rho,\,\sigma,\,\tau)$. - *
\hfill $19$--$2$

% ========================================================================
% PDF p. 168 = book p. 148
% ========================================================================
\medskip\noindent\textit{[book p.~148]}\quad \textbf{72.}\quad Collecting the terms, we have
\[
\partial_1\,[013] - \partial_1\,[103] =
\]
\[
\frac{1}{\Delta!}\left\{\sigma\,\Bigl[\rho\,(-\,8hn + 5jm + 2l^{2} - 2qr) + \frac{1}{q}\,(4phl - 3qj r - 3lhj + 2jh) + \frac{1}{q^{2}}\,(-\,2jhk + 4qhjn - 2jnl)\Bigr]\right.
\]
\[
+ \sigma\,\left[\frac{1}{q}\,(-\,4fp + 4jq - 2lm + nr) + \frac{1}{q}\,(2fhi - 3jp - 9jp + 6jpm + 2lr) + \frac{1}{q^{2}}\,(2k l m^{2} - 8 i k - 18np - 2lhr)\right]
\]
\[
+ \sigma\,\left[\frac{1}{q}\,((gr + 2hi + jp + n + 14hi) + \frac{1}{q}\,(-\,2pkn + 2lkj n - 2pnm) + \frac{1}{q^{2}}\,(2pkl - 2qhij - 2qjm + 2j^{2}n)\right]
\]
\[
\left.- *\right\}
\]

\medskip\noindent\textbf{73.}\quad The right-hand side depends on the points $0$, $1$, $3$ and $2$: viz.\ we have therein $\rho = 023$, $\Delta = 123$, \&c., but the left-hand side depending on only the points $0$, $1$ and $2$, the right-hand side cannot really contain $3$, and it must then remain unaltered, if for $2$ we substitute as another function of $0$, $1$, $3$ and $5$, viz.\ $\rho$, $\Delta$, $f$, \&c., will become $063$, $163$, $0/3$, \&c.; we have then $\phi(0,\,1,\,3) = $ the above linear function with $2$ thus replaced by $6$; say
\[
\phi(0,\,1,\,3) = \frac{1}{\Delta!}\,[\rho\,(\;) + \sigma(\;) + \tau(\;)],
\]
a given function of the points $0$, $1$, $3$ and the arbitrary point $6$, on the quartic curve; we therefore write it $\phi(0\,;\,1,\,3,\,6)$. There is no obvious value for $X(0,\,1,\,3)$ which will produce any simplification: I therefore take this function to be $= 0$; and the final result is
\[
Q_{k,\,01} = [012] + \Theta_{k,\,12} = -([025] + [031] + [012]),
\]
where $\Theta_{k,\,12}$ is a function integral and linear as regards the coordinates $(x,\,y,\,z)$ of the point $0$, but transcendental as regards the parametric points $1$, $2$; and containing besides the arbitrary points $3$, $6$, of the quartic curve, its value being determined by the differential formulae
\[
\partial_1\Theta_{k,\,12} = \phi(0',\,1,\,3,\,6),\quad \partial_2\Theta_{k,\,12} = -\,\phi(0',\,2,\,3,\,6),
\]
where $\phi(0',\,1,\,3,\,6)$ is a given function as above. I do not see the meaning of the very complicated linear function of $(\rho,\,\sigma,\,\tau)$, nor how to reduce it to any form such as the simple one $\partial_1\,[036]$, which presents itself in the case of the cubic curve.

\medskip
\begin{flushright}
\textit{Cambridge, England, October 5, 1882.}
\end{flushright}

% ========================================================================
% PDF p. 169 = book p. 149 -- Chapter IV. The Major Function (continued)
% ========================================================================
\subsection*{Chapter IV. The Major Function $(x,\,y,\,z)^{n-3}_k$ continued.}
\addcontentsline{toc}{subsection}{\quad Chapter IV. The Major Function (continued)}

\medskip\noindent\textit{The Conversion, Fixed Curve a Quartic, continued.}\quad Art.\ Nos.\ 74 to 82.

\medskip\noindent\textit{[book p.~149]}\quad \textbf{74.}\quad I resume the question considered \textit{ante} Nos.\ 66 to 73. The general problem, where the fixed curve is any given curve whatever, has recently been solved in a very complete and elegant form by Dr Nöther, in the two notes ``Zur Reduction algebraischer Differentialausdrücke auf die Normalformen'' and ``Ueber die algebraischen Differentialausdrücke, $2^{e}$ Note,'' \textit{Sitzungsb.\ der phys.-med.\ Soc.\ zu Erlangen}, 10 Dec.\ 1883 and 14 Jan.\ 1884. I consider here the case of the quartic curve, $n = 4$, and connect his result with my former investigations.

We have the differential
\[
\frac{(x,\,y,\,z)^{n-3}_k\,d\omega}{012} = \frac{\Omega_{12}\,d\omega}{012},
\]
where $\Omega_{12}$, or as I also write it $\Omega(0;\;1,\;2;\;3,\;4,\;5)$, is a rational and integral function of the degree $(n - 2) = 2$ in the current coordinates $(x,\,y,\,z)$: it depends also on the parametric points $1$, $2$, which are points on the quartic, coordinates $(x_1,\,y_1,\,z_1)$, $(x_2,\,y_2,\,z_2)$ respectively; and on $(p = )\,3$ other points $3$, $4$, $5$ are arbitrary points which we take to be in fact the points $(x_3,\,y_3,\,z_3)$, $(x_4,\,y_4,\,z_4)$, $(x_5,\,y_5,\,z_5)$ respectively. The curve $\Omega = 0$ is a conic, passing through the points $3$, $4$, $5$ and through the $(n - 2 =)\,2$ residues of the parametric points 1, 2, of the quartic curve, $n = 4$, and connect his result with my former investigations.

We have hence
\[
\Omega\,(0;\;1,\;2;\;3,\;4,\;5) = \frac{1}{\Delta}\,(a\,1^{n-2}\,2\,1\,1\,1,\;\ldots\,)\,(x,\,y,\,z)^{n-2};\;\text{see No.\ 2.}
\]
the function determined by the equation
\[
\left|\;\begin{array}{cc}
(x,\,y,\,z)^{2} & ,\quad \Omega\\
1\,(x_1,\,y_1,\,z_1)^{2} & ,\quad 4\,.\,1^{3}\,2\\
2\,(x_1,\,y_1,\,z_1\;\xi\,x_2,\,y_2,\,z_2),\;6\,.\,1^{2}\,2^{2}\\
1\,(x_2,\,y_2,\,z_2)^{2} & ,\quad 4\,.\,1\,2^{3}\\
(x_3,\,y_3,\,z_3)^{2} & ,\quad 0\\
(x_4,\,y_4,\,z_4)^{2} & ,\quad 0\\
(x_5,\,y_5,\,z_5)^{2} & ,\quad 0
\end{array}\;\right| = 0,
\]
\bigskip
\hrule
\smallskip
\noindent\textsuperscript{*}\,$\bigl(x_a\,\dfrac{d}{dx_b} + y_a\,\dfrac{d}{dy_b} + z_a\,\dfrac{d}{dz_b}\bigr)(x_b,\,y_b,\,z_b)^{n-2} = $ see No.\ 2.

% ========================================================================
% PDF p. 170 = book p. 150
% ========================================================================
\medskip\noindent\textit{[book p.~150]}\quad this is of the form
\[
M\Omega + \square = 0,
\]
and as appears in No.\ 46, $M = 123\,.\,124\,.\,125\,.\,345$. Hence writing $\square\,(0;\;1,\;2;\;3,\;4,\;5)$ for $\square$, we have
\[
\Omega_n = \Omega\,(0;\;1,\;2;\;3,\;4,\;5) = \frac{-\,\square\,(0;\;1,\;2;\;3,\;4,\;5)}{123\,.\,124\,.\,125\,.\,345}.
\]
Hence further writing
\[
Q_n = Q\,(0;\;1,\;2;\;3,\;4,\;5) = \frac{\Omega\,(0;\;1,\;2;\;3,\;4,\;5)}{012},
\]
so that the differential is $Qd\omega$, $= Q\,(0;\;1,\;2;\;3,\;4,\;5)\,d\omega$, we have
\[
Q = Q\,(0;\;1,\;2;\;3,\;4,\;5) = \frac{-\,\square\,(0;\;1,\;2;\;3,\;4,\;5)}{012\,.\,123\,.\,124\,.\,125\,.\,345},
\]
which is of the form
\[
Q = \frac{012\bdet{2345}}{012\bdet{2345}},\;= 012\bdet{2345},
\]
viz.\ $Q$ is a rational fraction where the numerator is of the degree $2$, and the denominator of the degree $5$ as regards the coordinates $(x,\,y,\,z)$ of the current point; but the numerator and denominator are each of the degree $4$ as regards the coordinates of the points $1$, $2$ separately; and the degree $2$ as regards the coordinates of the points $3$, $4$, $5$ separately.

\medskip\noindent\textbf{76.}\quad The signification of the symbol of quasi-differentiation $\partial$ (applicable only to a function of the degree $0$ in the coordinates to which the differentiations have reference) is explained \textit{ante} No.\ 60. The function $Q$ just mentioned is of the degree $0$ in regard to the coordinates of each of the points $1$, $2$, $3$, $4$, $5$; and it can thus be operated upon by the symbols $\partial_1$, $\partial_2$, $\partial_3$, $\partial_4$, $\partial_5$ respectively. Observe, in particular, that we have $\partial_3 Q\,(0;\;1,\;2;\;3,\;4,\;5) = 012\bdet{2345}$, viz.\ it is of the degree $1$ in the coordinates of the points $0$ and $1$ respectively, but of the degree $0$ in regard to the coordinates of the points $2$, $3$, $4$, $5$ respectively.

\medskip\noindent\textbf{77.}\quad This being so, we may consider the function
\[
H\,(0;\;1,\;2;\;3,\;4,\;5;\;6,\;7,\;8) = \partial_3 Q\,(0;\;1,\;2;\;3,\;4,\;5)
\]
\[
+ \partial_3 Q\,(1;\;3,\;2;\;6,\;4,\;5)\,.\,\frac{045}{345}
\]
\[
+ \partial_4 Q\,(1;\;4,\;2;\;3,\;7,\;5)\,.\,\frac{053}{453}
\]
\[
+ \partial_5 Q\,(1;\;5,\;2;\;3,\;4,\;8)\,.\,\frac{034}{534},
\]
where $6$, $7$, $8$ are arbitrary points on the quartic; the functions
\[
\partial_3 Q\,(1;\;3,\;2;\;6,\;4,\;5),\;\partial_4 Q\,(1;\;4,\;2;\;3,\;7,\;5),\;\partial_5 Q\,(1;\;5,\;2;\;3,\;4,\;8),
\]

% ========================================================================
% PDF p. 171 = book p. 151
% ========================================================================
\medskip\noindent\textit{[book p.~151]}\quad are functions of the same form as $\partial_3 Q\,(0;\;1,\;2;\;3,\;4,\;5)$, and derived from it by changing in each case the current point $0$ into the parametric point $1$, and the points $3$, $4$, $5$ respectively, and using the corresponding point $3$, $4$ or $5$ by the new arbitrary point $6$, $7$ or $8$. Further $045$, \&c., denote determinants as above; so that in $H$ each of the last three terms is, in fact, as regards the point $0$, a mere linear function of the coordinates $(x,\,y,\,z)$ of this point.

We have $\partial_1 Q\,(1;\;3,\;2;\;6,\;4,\;5) = 13'\bdet{2456}$, and hence this function multiplied by $\dfrac{045}{345}$ is $= 01'\bdet{2456}$; and so for the third and fourth terms of $H$: thus each of the four terms of $H$ is $= 01'\bdet{2456789}$, of the degree $1$ in the coordinates of the points $0$ and $1$ respectively, but of the degree $0$ in regard to the coordinates of the other points $2$, $3$, $4$, $5$, $6$, $7$, $8$.

Nöther's conversion-theorem consists herein, that the function
\[
H\,(0;\;1,\;2;\;3,\;4,\;5;\;6,\;7,\;8)
\]
is unaltered by the interchange of the two points $0$, $1$; or putting for shortness
\[
H\,(0;\;1,\;2;\;3,\;4,\;5;\;6,\;7,\;8) = H_1\,(0),
\]
the theorem is
\[
H_1\,(0) = H_0\,(1).
\]

\medskip\noindent\textbf{78.}\quad We have, No.\ 59,
\[
\frac{\tfrac{1}{2}\Omega_{12}}{012} = \frac{-\,01^{3}\,.\,02^{2} + 01^{2}\,2\,.\,012 + 01\,2\,.\,1^{2}\,2}{012\,.\,1^{2}\,2},\;= [012],
\]
or as for greater simplicity I write it
\[
= 012,
\]
viz.\ $012$ is now written instead of $[012]$ to denote the function just given as the value of $\dfrac{\tfrac{1}{2}\Omega_{12}}{012}$. $\Omega_{12}$ is thus $= r\,.\,012\,.\,012$, viz.\ this is a particular form of $\Omega_{12}$ satisfying the conditions that $\Omega_{12} = 0$ is a conic passing through the residues of the points $1$, $2$ on writing therein $(x_1,\,y_1,\,z_1)$ for $(x,\,y,\,z)$ becomes $= s\,.\,012\,1^{2}\,2 +$ arbitrary linear function $(x,\,y,\,z)$. The before-mentioned function $\Omega\,(0;\;1,\;2;\;3,\;4,\;5)$ is a function satisfying these conditions on writing therein for $4$, $5$, $6$ in succession the values $3$, $4$, $5$, becomes $= 0$; we have therefore
\[
\frac{\tfrac{1}{2}\Omega\,(0;\;1,\;2;\;3,\;4,\;5)}{012} = 012 - 012\,\frac{034}{345} - 012\,\frac{054}{354} - 012\,\frac{045}{354},
\]
viz.\ the value of $Q$ given by this equation, on writing therein $0 = 3$, $4$, or $5$, becomes $= 0$ as it should do.

% ========================================================================
% PDF p. 172 = book p. 152
% ========================================================================
\medskip\noindent\textit{[book p.~152]}\quad \textbf{79.}\quad We thus have
\[
Q\,(0;\;1,\;2;\;3,\;4,\;5) = 012 - 012\,\frac{045}{345} - 012\,\frac{053}{453} - 012\,\frac{034}{534},
\]
and Nöther's conversion-equation becomes
\[
\partial_1\,\bigl\{012 - 012\,\frac{045}{345} - 012\,\frac{053}{453} - 012\,\frac{034}{534}\bigr\}
\]
\[
+ \partial_3\,\Bigl\{1'02 - 6'32\,\frac{145}{645} - 4'32\,\frac{156}{456} - 5'32\,\frac{164}{564}\Bigr\}\frac{045}{345}
\]
\[
+ \partial_4\,\bigl\{1'02 - 3'42\,\frac{175}{375} - 7'42\,\frac{153}{753} - 5'42\,\frac{137}{573}\bigr\}\frac{053}{453}
\]
\[
+ \partial_5\,\bigl\{1'02 - 3'52\,\frac{148}{348} - 4'52\,\frac{183}{483} - 8'52\,\frac{134}{834}\bigr\}\frac{034}{534}
\]
\[
=\;\partial_0\,\bigl\{1'02 - 6'32\,\frac{145}{645} - 4'32\,\frac{156}{456} - 5'32\,\frac{164}{564}\bigr\}
\]
\[
+ \partial_3\,\bigl\{0'62 - 6'32\,\frac{045}{645} - 4'32\,\frac{056}{456} - 5'32\,\frac{064}{564}\bigr\}\frac{045}{345}
\]
\[
+ \partial_4\,\bigl\{1'02 - 3'72\,\frac{075}{375} - 7'42\,\frac{053}{753} - 5'72\,\frac{037}{573}\bigr\}\frac{053}{453}
\]
\[
+ \partial_5\,\bigl\{0'82 - 3'82\,\frac{048}{348} - 4'82\,\frac{083}{483} - 8'82\,\frac{034}{834}\bigr\}\frac{034}{534},
\]
an equation where the functions operated on with the $\partial$'s are only functions such as $012$; for there is not any determinant operated upon containing the number which is the suffix of the $\partial$ operating upon it.

\medskip\noindent\textbf{80.}\quad Taking all the terms over to the left-hand side, there are in all $32$ terms; but of these $3 + 3$ destroy each other, and $6 + 6$ unite in pairs into $6$ terms: there are thus in all $7 + 7 + 6$, $= 20$ terms; viz.\ multiplying the whole equation by $345$, it is found that the equation becomes:
\[
\text{as so this may be written}\qquad\text{where}
\]
\[
345\,(\;\partial_1\,012 - \partial_0\,1'02\bigr)\quad\;\;\;345\;\;\boxed{012}\;\;\;\;\;\;\;= \partial_1\,012 - \partial_0\,1'02,\,\&\text{c.}
\]
\[
-\,045\,(-\,\partial_1\,102 + \partial_0\,012)\quad -\,045\;\;\boxed{312}
\]
\[
-\,053\,(-\,\partial_1\,142 + \partial_0\,012)\quad -\,305\;\;\boxed{412}
\]
\[
-\,034\,(-\,\partial_1\,152 + \partial_0\,012)\quad -\,340\;\;\boxed{512}
\]

% ========================================================================
% PDF p. 173 = book p. 153
% ========================================================================
\medskip\noindent\textit{[book p.~153]}\quad
\[
-\,145\,(\;\partial_1\,032 - \partial_2\,302)\quad -\,145\;\;\boxed{032}
\]
\[
-\,153\,(\;\partial_1\,042 - \partial_2\,402)\quad -\,315\;\;\boxed{042}
\]
\[
-\,134\,(\;\partial_1\,052 - \partial_2\,502)\quad -\,341\;\;\boxed{052}
\]
\[
+\,013\,(-\,\partial_1\,542 + \partial_2\,452)\quad +\,301\;\;\boxed{452}
\]
\[
+\,014\,(-\,\partial_1\,352 + \partial_2\,532)\quad +\,140\;\;\boxed{532}
\]
\[
+\,015\,(-\,\partial_1\,432 + \partial_2\,342)\quad +\,015\;\;\boxed{342}
\]
\[
= 0,\hspace{8em}= 0,
\]
viz.\ the equation is
\[
\Sigma \pm 345\;\;\boxed{012}\;\; = 0.
\]
the nine terms which follow the first term $345\;\boxed{012}$ of the sum being obtained by the interchanges of $0$, $1$ (one or each) with the $3$, $4$, $5$, each interchange giving rise to a sign $-$.

\medskip\noindent\textbf{81.}\quad In obtaining the foregoing result, we have, for instance, a pair of terms
\[
\partial_1\,542\,\frac{-\,137\,.\,035 + 037\,.\,135}{537},\;= \partial_1\,542\,\frac{-013\,.\,537}{537},\;= \partial_1\,542\,(-\,013),
\]
viz.\ this depends on the equation
\[
-\,137\,.\,035 + 037\,.\,135 - 537\,.\,013 = 0,
\]
or say
\[
-\,137\,.\,035 + 037\,.\,135 + 013\,.\,357 = 0,
\]
an identity which, in a form which will be readily understood, may be written
\[
\det\,\left|\;\begin{array}{c}0137\\013735\end{array}\;\right| = 0.
\]
Similarly, the two terms which contain $\partial_3\,452$ combine into the single term $\partial_3\,452\,(013)$; and the two new terms taken together are
\[
013\,(-\,\partial_3\,542 + \partial_2\,452) = 301\;\;\boxed{452}.
\]
C.\ XII. \hfill 20

% ========================================================================
% PDF p. 174 = book p. 154
% ========================================================================
\medskip\noindent\textit{[book p.~154]}\quad \textbf{82.}\quad The proof of the identity,
\[
\Sigma \pm 345\;\;\boxed{012}\;\; = 0,
\]
depends on the property of the function
\[
\boxed{012}_{1},\;= \partial_1 012 - \partial_0 1'02,
\]
enunciated No.\ 67, and proved \textit{à posteriori} by the tedious calculation Nos.\ 69 to 73, viz.\ in No.\ 67, writing $2$ in place of $3$, this is---$\partial_1\,012 - \partial_0\,1'02$ is equal to the difference of two functions, the first of them linear in the coordinates $(x,\,y,\,z)$ of the point $0$, but depending also on the coordinates of the points $1$ and $2$, the second of them linear in the coordinates $(x_1,\,y_1,\,z_1)$ of the point $1$, but depending also on the coordinates of the points $0$ and $2$. Or, what is the same thing, the property is
\[
\boxed{012}_{1},\;= A_{12}x + B_{12}y + C_{12}z - (A_{01}x_1 + B_{01}y_1 + C_{01}z_1),
\]
where $A_{12}$, $B_{12}$, $C_{12}$ are functions of $(x_1,\,y_1,\,z_1)$, $(x_2,\,y_2,\,z_2)$, and $A_{01}$, $B_{01}$, $C_{01}$ are the like functions of $(x,\,y,\,z)$, $(x_2,\,y_2,\,z_2)$.

Substituting such values in the sum $\Sigma \pm 345\;\boxed{012}$, but writing down only the terms which contain $x$, these are
\[
345\,(A_{12}x - A_{01}x_1)
\]
\[
-\,045\,(A_{12}x_1 - A_{01}x_1)\qquad -\,145\,(A_{02}x - A_{20}x_0)\qquad +\,301\,(A_{42}x_4 - A_{42}x_5)
\]
\[
-\,305\,(A_{42}x_1 - A_{01}x_4)\qquad -\,315\,(A_{02}x - A_{40}x_0)\qquad +\,140\,(A_{52}x_5 - A_{52}x_3)
\]
\[
-\,340\,(A_{52}x_1 - A_{01}x_5)\qquad -\,341\,(A_{02}x - A_{50}x_0)\qquad +\,015\,(A_{32}x_3 - A_{32}x_4),
\]
This is
\[
=\quad A_{12}\,(\,x\,.\,345 - x_1\,.\,015 + x_2\,.\,015)
\]
\[
+\;A_{02}\,(\;x\,.\,345 - x_0\,.\,345 - x_2\,.\,340)
\]
\[
+\;A_{42}\,(\,x_4\,.\,145 - x_1\,.\,305 + x_0\,.\,340)
\]
\[
+\;A_{52}\,(\;x_5\,.\,134 + x_1\,.\,340 + x_0\,.\,140)
\]
\[
+\;A_{01}\,(-\,x_1\,.\,345 + x_4\,.\,305 + x_5\,.\,340)
\]
\[
+\;A_{42}\,(-\,x\,.\,341 + x_5\,.\,140 + x_3\,.\,015),
\]
where the coefficient of each of the $A$'s is identically $= 0$; and similarly, the terms in $y$ and the terms in $z$ are each $= 0$. We have thus the proof of the identity
\[
\Sigma \pm 345\;\;\boxed{012}\;\; = 0,
\]
that is, of the conversion-equation $H_1(0) = H_0(1)$.

% ========================================================================
% PDF p. 175 = book p. 155
% ========================================================================
\medskip\noindent\textit{[book p.~155]}\quad \textit{The Syzygy---Fixed Curve a Quartic.}\quad Art.\ No.\ 83.

\medskip\noindent\textbf{83.}\quad I revert to the theory of the Syzygy, \textit{ante} No.\ 59.

We have
\[
Q\,(0;\;1,\;2;\;3,\;4,\;5) = 012 - 012\,\frac{045}{345} - 012\,\frac{053}{453} - 012\,\frac{034}{534},
\]
or if for convenience we take instead of $1$, $2$, the parametric points to be $a$, $\beta$ coordinates $(x_a,\,y_a,\,z_a)$ and $(x_\beta,\,y_\beta,\,z_\beta)$ respectively, then this equation is
\[
Q_{a\beta} = Q\,(0;\;a,\;\beta;\;3,\;4,\;5) = 0a\beta - 0a\beta\,\frac{045}{345} - 0a\beta\,\frac{053}{453} - 0a\beta\,\frac{034}{534}.
\]
Considering a new parametric point $\gamma$, and forming the like functions $Q_{\beta\gamma}$ and $Q_{\gamma a}$, it is to be shown that we have identically
\[
Q_{a\beta} + Q_{\beta\gamma} + Q_{\gamma a} = 0.
\]
To prove this, observe that, in the equation at the end of No.\ 59, $\Delta$, $\rho$, $\sigma$, $\tau$ denote $123$, $023$, $031$, $012$ respectively. Hence writing therein $a$, $\beta$, $\gamma$ in place of $1$, $2$, $3$ respectively, and putting $A$, $B$, $C$ for the coefficients (including therein the factor $\dfrac{1}{\Delta!}$) of $\rho$, $\sigma$, $\tau$ respectively, the equation is
\[
0\bar{\beta}\gamma + 0\gamma a + 0a\beta = A\,.\,0\beta\gamma + B\,.\,0\gamma a + C\,.\,0a\beta;
\]
where $A$, $B$, $C$ are absolute constants (functions, that is, of the coefficients of the quartic) each divided by $(a\beta\gamma)^{6}$. We hence obtain
\[
(Q_{\beta\gamma} + Q_{\gamma a} + Q_{a\beta})\,.\,345 = 345\,(A\,.\,0\beta\gamma + B\,.\,0\gamma a + C\,.\,0a\beta)
\]
\[
-\,045\,(A\,.\,3\beta\gamma + B\,.\,3\gamma a + C\,.\,3a\beta)
\]
\[
-\,053\,(A\,.\,4\beta\gamma + B\,.\,4\gamma a + C\,.\,4a\beta)
\]
\[
-\,034\,(A\,.\,5\beta\gamma + B\,.\,5\gamma a + C\,.\,5a\beta).
\]
On the left-hand side the whole coefficient of $A$ is $= 0$; viz.\ the coefficient has the value $\det\,\left|\begin{array}{c}0345\\0345\beta\gamma\end{array}\right|$, which is $= 0$. Similarly, the whole coefficient of $B$ is $= 0$, and the whole coefficient of $C$ is $= 0$; and we have thus the required result
\[
Q_{\beta\gamma} + Q_{\gamma a} + Q_{a\beta} = 0.
\]
The syzygy is thus obtained in a more perfect form than in No.\ 59; viz.\ by considering (instead of $0\bar{a}\beta$) the new form $Q_{a\beta}$ then, instead of a sum which is a linear function of the coordinates $(x,\,y,\,z)$, we obtain a sum $= 0$.
\hfill $20$--$2$

\end{document}
